% answers_short.tex —— 2006 年·普通高等学校招生考试·浙江卷理科 · 简短答案速查
% 由 tools/make_short_answers.py 自 answers.tex 的 \gksolhead/\gkanswer 提取，
% 并与 2006/2006_数学真题_浙江理_答案要点.md 逐题比对后生成，勿手改（改源 MD 后重跑）。

\begin{gknotice}
\textbf{使用说明}：本节只给最终答案便于快速核对，完整推导见解析版。答案为编写者独立推导，\textbf{非官方标准答案}。
\end{gknotice}

\gkpart{选择题}
\begin{center}\begin{tabular}{ccccccccccc}
\toprule
题号 & 1 & 2 & 3 & 4 & 5 & 6 & 7 & 8 & 9 & 10 \\
\midrule
答案 & A & C & A & B & C & C & A & D & B & D \\
\bottomrule
\end{tabular}\end{center}

\gkpart{填空题}
\begin{itemize}[leftmargin=2.6em]
  \item[\textbf{11.}] $-1$。
  \item[\textbf{12.}] $3/2$。
  \item[\textbf{13.}] $4$。
  \item[\textbf{14.}] $[\sqrt2/4,1/2]$。
\end{itemize}

\gkpart{解答题}
\begin{itemize}[leftmargin=2.6em]
  \item[\textbf{15.}] (1) $\varphi=\pi/6$；(2) $\arccos(15/17)$
  \item[\textbf{16.}] (1) $a>0$ 且 $-2<b/a<-1$；(2) 方程 $f(x)=0$ 在 $(0,1)$ 内有两个实根
  \item[\textbf{17.}] (1) 见解析版；(2) $CD$ 与平面 $ADMN$ 所成的角为 $\arcsin\sqrt{2/5}$
  \item[\textbf{18.}] (1) $1/60$；(2) $n=2$
  \item[\textbf{19.}] (1) $\frac{x^2}{2}+2y^2=1$；(2) 见解析版
  \item[\textbf{20.}] (1) $x_n^2+x_n=3x_{n+1}^2+2x_{n+1}$；(2) $(1/2)^{n-1}\leqslant x_n\leqslant(1/2)^{n-2}$
\end{itemize}
